% Revize 2026-11-01: hodinove poradi 73--75 podle gradace obtiznosti.
% Zadani jsou beze zmen prevzata z predchozich hodin 74, 75; zadne nebylo ztraceno.
% Metodicka navaznost: hodina 72 (ekvivalentni upravy), Realisticky.cz 020203.
\documentclass[11pt,a4paper]{article}
\input{../../sablony/pripravy/priprava.sty}
\nastavPripravu{1. ročník}{Rovnice a nerovnice}{73}{Lineární rovnice I: základy}
\begin{document}
\blok{Návaznost: od ekvivalentních úprav k řešení}
V hodině 72 jsme zavedli rovnici, její kořen a ekvivalentní úpravy. Dnes pravidla použijeme na jednoduchých rovnicích. U každé úpravy má žák vysvětlit, \emph{co provádí s oběma stranami}; fráze „převedu na druhou stranu a změním znaménko“ není sama o sobě odůvodněním.
\begin{definiceBox}
\textbf{Lineární rovnice o jedné neznámé} je rovnice, kterou lze ekvivalentními úpravami převést na tvar $ax+b=0$, kde $a,b\in\mathbb R$. Při $a\ne0$ má právě jeden kořen; pro $a=0$, $b\ne0$ nemá žádný; pro $a=b=0$ vyhovuje každé číslo z přípustného oboru.
\end{definiceBox}
\blok{Tři možné výsledky}
Z tvaru $ax+b=0$ plyne: je-li $a\ne0$, řešením je $x=-b/a$; rovnost $0=5$ je spor; rovnost $0=0$ je pravdivá pro všechna přípustná $x$. Tento rozbor se objeví i u složitějších rovnic, nevyplatí se jej vynechat.
\begin{prikladReseniBox}[title={Příklad 1 – jeden kořen}]
Vyřešte $3x+7=22$.
\tcblower
\begin{align*}
3x+7&=22 &&\big|\,-7\\
3x&=15 &&\big|\,:3\\
x&=5.
\end{align*}
Zkouška: $3\cdot5+7=22$.
\end{prikladReseniBox}
\begin{prikladReseniBox}[title={Příklad 2 – závorka a neznámá na obou stranách}]
Vyřešte $4-2(x-1)=3x+1$.
\tcblower
\begin{align*}
4-2x+2&=3x+1\\
6-2x&=3x+1\\
5&=5x\\
x&=1.
\end{align*}
Zkouška: $4-2(1-1)=4=3\cdot1+1$. Záporný násobitel mění znaménko \emph{obou} členů v závorce.
\end{prikladReseniBox}
\poznamka{První příklad vyřeš krátce; většinu společné práce věnuj druhému a otázce, proč jsou úpravy ekvivalentní. Závorky jsou zde jen v jedné úrovni; složené závorky až v hodině 74.}
\blok{Samostatné procvičování – stupňované úlohy}
\begin{zadaniReseniBox}
\begin{enumerate}[label=\alph*)]
\item $2x-5=-9$.
\item $3(x+2)=2x+9$.
\item $5x+1=5x-2$.
\end{enumerate}
\tcblower
\begin{enumerate}[label=\alph*)]
\item $\begin{aligned}2x-5&=-9\\2x&=-4\\x&=-2,\quad K=\{-2\}.\end{aligned}$
\item $\begin{aligned}3x+6&=2x+9\\x+6&=9\\x&=3,\quad K=\{3\}.\end{aligned}$
\item $\begin{aligned}5x+1&=5x-2\\1&=-2\\K&=\varnothing.\end{aligned}$
\end{enumerate}
\end{zadaniReseniBox}
\blok{Volitelná rezerva – rovnice jako identita}
\begin{zadaniReseniBox}
Pro jaké $x$ platí $2(x-1)+3=2x+1$?
\tcblower
\begin{align*}
2x-2+3&=2x+1\\
2x+1&=2x+1\\
0&=0\\
K&=\mathbb R.
\end{align*}
Na rozdíl od sporu je poslední rovnost pravdivá vždy.
\end{zadaniReseniBox}
\blok{Domácí práce – základní rovnice a počet řešení}
\begin{zadaniReseniBox}
\begin{enumerate}[label=\alph*)]
\item $7x-4=3x+8$.
\item $3(x-2)=2x+5$.
\item $5-2(3x-4)=3(x-2)$.
\item $4x+2=4x+2$.
\item $2(x-3)-3(x+1)=5-x$.
\end{enumerate}
\tcblower
\begin{enumerate}[label=\alph*)]
\item $\begin{aligned}7x-4&=3x+8\\4x&=12\\x&=3.\end{aligned}$
\item $\begin{aligned}3x-6&=2x+5\\x&=11.\end{aligned}$
\item $\begin{aligned}13-6x&=3x-6\\19&=9x\\x&=\tfrac{19}{9}.\end{aligned}$
\item $\begin{aligned}4x+2&=4x+2\\0&=0\\K&=\mathbb R.\end{aligned}$
\item $\begin{aligned}2x-6-3x-3&=5-x\\-x-9&=5-x\\-9&=5\\K&=\varnothing.\end{aligned}$
\end{enumerate}
\end{zadaniReseniBox}
\poznamka{45 minut: úvod a rozbor $ax+b=0$ 8 min, dva řešené příklady 13 min, samostatně 15 min, diskuse o sporu a zkoušce 6 min, zadání domácí práce 3 min. Rezerva je navíc. Následuje hodina 74: postupné odstraňování vnořených závorek.}
\end{document}
